% example-2026-1-A.tex — 还原 2025-2026 学年第 1 学期概率论与数理统计 A 卷
\documentclass[answerspace=14]{shnu-exam}

\examtitle{概率论与数理统计}
\examtype{A卷}
\semester{2025~2026 学年~~第~1~学期}
\examdate{2026年1月13日}
\examduration{90}

\begin{document}
\maketitle

\section{一、选择题（每小题2分，共16分）}

\question 以下说法中正确的是
\choice{概率为0的事件是不可能事件}
\choice{相互独立的事件不可能同时发生}
\choice{泊松分布的样本点数量是无限的}
\choice{以上均不对}

\question 以下关于指数分布的描述中错误的是
\choice{指数分布具有无记忆性}
\choice{指数分布的数学期望与方差可能相同}
\choice{密度函数为0的点概率为0}
\choice{参数$\theta$表示事件发生的平均次数}

\question 二维随机变量$(X,Y)$的联合概率密度函数为
$f(x,y)=\begin{cases}2,&0<x<y<1\\0,&\text{其他}\end{cases}$
下列结论中正确的是
\choice{$X$与$Y$相互独立}
\choice{$P\{X+Y<1\}=\frac12$}
\choice{$P\{X<\frac12,Y<\frac12\}=\frac12$}
\choice{$P\{X<Y\}=1$}

\question 设随机变量$X$的概率密度函数为
$f(x)=\begin{cases}2x,&0<x<1\\0,&\text{其他}\end{cases}$
下列结论中正确的是
\choice{$E(X)=\frac12$}
\choice{$E(X^2)=\frac12$}
\choice{$D(X)=\frac{1}{18}$}
\choice{$D(X)=\frac{1}{12}$}

\question 设随机变量$X_1,X_2,\dots,X_n$相互独立同分布，且$E(X_i)=\mu$，$D(X_i)=\sigma^2$，$\bar{X}=\frac1n\sum X_i$。当$n$足够大时，下列结论中正确的是
\choice{$\bar{X}$近似服从$N(\mu,\sigma^2)$}
\choice{$\sqrt{n}(\bar{X}-\mu)$近似服从$N(0,\sigma^2)$}
\choice{$\bar{X}-\mu$近似服从$N(0,\frac{\sigma^2}{n})$}
\choice{$\sum X_i$近似服从$N(n\mu,n^2\sigma^2)$}

\question 设随机样本$X_1,X_2,\dots,X_n$来自正态总体$N(\mu,\sigma^2)$，其中$\mu,\sigma^2$均未知。设$\bar{X}=\frac1n\sum X_i$，$S^2=\frac1{n-1}\sum(X_i-\bar{X})^2$。下列统计量中，服从$F$分布的是
\choice{$\frac{(n-1)\bar{X}}{\sigma^2}$}
\choice{$\frac{\bar{X}-\mu}{S/\sqrt{n}}$}
\choice{$\frac{S_1^2/\sigma_1^2}{S_2^2/\sigma_2^2}$}
\choice{$\frac{S_1^2}{S_2^2}$}

\question 设随机样本$X_1,X_2,\dots,X_n$来自正态总体$N(\mu,\sigma^2)$，其中$\sigma^2$已知，$\mu$未知。下列关于参数$\mu$的估计量中，既是无偏估计，又具有最小方差的是
\choice{$X_1$}
\choice{$\bar{X}$}
\choice{$\frac1n\sum|X_i|$}
\choice{$\frac{X_1+X_n}{2}$}

\question 已知某工厂生产的零件直径服从正态分布，已知总体方差为$\sigma^2=4$。现从该工厂随机抽取$n=25$个零件，测得样本均值为$\bar{x}=50.7$。现要检验该工厂零件的平均直径是否为50，显著性水平取$\alpha=0.05$，上$0.05$分位点为$1.634$，上$0.025$分位点为$1.96$。设原假设$H_0:\mu=50$，备择假设$H_1:\mu\neq50$，下列结论中正确的是
\choice{接受原假设$H_0$}
\choice{拒绝原假设$H_0$}
\choice{无法进行假设检验}
\choice{需要增加样本容量后才能判断}

\section{二、填空题（每空3分，共24分）}

\question 已知$P(A)=0.6$，$P(B)=0.5$，$P(AB)=0.2$。则$P(A\bar{B})=$\fillin.

\question 设连续型随机变量$X$的概率密度函数为
$f(x)=\begin{cases}k(2-x),&0<x<2\\k(x-2),&2\leq x<4\\0,&\text{其他}\end{cases}$
则常数$k$为\fillin.

\question 设二维随机变量$(X,Y)$的联合概率密度函数为
$f(x,y)=\begin{cases}2,&0<x<y<1\\0,&\text{其他}\end{cases}$
则当$0<x<y<1$时，$F(x,y)=$\fillin.

\question 设连续型随机变量$X$的概率密度函数为
$f(x)=\begin{cases}2(1-x),&0<x<1\\0,&\text{其他}\end{cases}$
记$Y=\min\{X,1-X\}$，则$E(Y)=$\fillin.

\question 设随机变量$X_1,X_2,\dots,X_n$相互独立同分布，且$E(X_i)=\mu$，$D(X_i)=\sigma^2$，根据中心极限定理，当$n$充分大时，随机变量$\frac{X_1+X_2+\cdots+X_n-n\mu}{\sqrt{n}\sigma}$近似服从\fillin 分布.

\question 设随机变量$X_1,X_2,\dots,X_n$相互独立，且都服从正态分布$N(\mu,\sigma^2)$，其中$\mu$已知，$\sigma^2$未知。则统计量$\sum\frac{(X_i-\mu)^2}{\sigma^2}$服从的分布是\fillin.

\question 设总体$X$服从指数分布，其中参数$\theta$未知。从该总体中抽取容量为$n$的简单随机样本$X_1,X_2,\dots,X_n$，则参数$\theta$的最大似然估计为\fillin.

\question 设对总体均值$\mu$进行假设检验，原假设与备择假设分别为$H_0:\mu=\mu_0$和$H_1:\mu\neq\mu_0$，显著性水平为$\alpha$，则该检验属于\fillin 检验.

\section{三、计算题（本题10分）}

\problem[lines=12]{某实验室有三台检测设备A、B、C。在所有检测任务中，有50\%由设备A完成，30\%由设备B完成，剩下20\%由设备C完成。已知设备A的误判率为1\%，设备B的误判率为2\%，设备C的误判率为4\%。现从所有已经完成的检测任务中随机抽取一次，发现该次检测结果是误判。请问这次检测是由设备B完成的可能性有多大？}

\section{四、计算题（本题16分）}

\problem[lines=18]{设二维随机变量$(X,Y)$的联合概率密度函数为
$f(x,y)=\begin{cases}4xy,&0<x<1,x^2<y<x\\0,&\text{其他}\end{cases}$
(1) 求$X$和$Y$的边缘密度函数；
(2) $X$和$Y$是否独立？}

\section{五、计算题（本题16分）}

\problem[lines=18]{设随机变量$X_1,X_2,\dots,X_n$独立同分布，其概率密度函数为
$f(x;\theta)=\begin{cases}\theta x^{\theta-1},&0<x<1\\0,&\text{其他}\end{cases}$
分别用矩估计和最大似然估计求参数$\theta$。}

\section{六、计算题（本题18分）}

\problem[lines=20]{设二维随机变量$(X,Y)$的联合概率密度函数为：
$f(x,y)=\begin{cases}c(x^2+y^2),&0<x<1,0<y<1,x+y\leq1\\0,&\text{其他}\end{cases}$
(1) 求常数$c$；
(2) 求数学期望$E[Y\mid X=x]$。}

\end{document}
